\chapter{Optimisation}

L'optimisation est au cœur de l'apprentissage automatique : on cherche en permanence à minimiser une fonction de coût (ou perte) par rapport aux paramètres d'un modèle.

\section{Convexité}

\begin{definition}[Fonction Convexe]
Une fonction $f : \mathbb{R}^n \rightarrow \mathbb{R}$ est convexe si, pour tous $x, y \in \mathbb{R}^n$ et pour tout $\theta \in [0, 1]$, on a :
$$ f(\theta x + (1-\theta)y) \leq \theta f(x) + (1-\theta)f(y) $$
\end{definition}

\begin{theorem}[Minimum Global]
Si une fonction de coût $J$ est strictement convexe, alors tout minimum local de $J$ est un minimum global unique.
\end{theorem}

\section{Descente de Gradient}

La descente de gradient est l'algorithme d'optimisation du premier ordre le plus utilisé pour minimiser une fonction de coût différentiable $J(\theta)$.

\begin{definition}[Mise à jour du gradient]
L'itération de la descente de gradient est donnée par :
$$ \theta_{t+1} = \theta_t - \alpha \nabla J(\theta_t) $$
où $\alpha$ est le taux d'apprentissage (learning rate) et $\nabla J(\theta_t)$ est le gradient de la fonction de coût au point $\theta_t$.
\end{definition}

\section{Multiplicateurs de Lagrange}

Pour les problèmes d'optimisation sous contraintes, on utilise la méthode des multiplicateurs de Lagrange.

\begin{definition}[Lagrangien]
Pour minimiser $f(x)$ sous la contrainte d'égalité $g(x) = 0$, on introduit le Lagrangien :
$$ \mathcal{L}(x, \lambda) = f(x) + \lambda g(x) $$
Les points critiques sont trouvés en résolvant $\nabla \mathcal{L}(x, \lambda) = 0$.
\end{definition}
